\documentclass[11pt]{article}
\usepackage{problem_set_style}
\usepackage{kotex}
\usepackage{booktabs}
\usepackage{array}
\begin{document}
\heading{G2}{중간 · 기말 공통}{}{Spring 2027}{리뷰(동료평가) 가이드라인}

\noindent{\small 각자 다른 팀의 논문 2편을 KSA Review에서 심사한다. 이 문서는 The Web Conference 2027 심사 지침, ICLR 2026$\cdot$2027, ICML 2026, NeurIPS 2026의 AI 정책을 우리 수업에 맞게 고친 것이다 (끝의 표 참고).}

\section*{먼저 이것만은}
\begin{enumerate}
\item \textbf{여러분의 리뷰가 그 팀이 받는 피드백의 전부다.} 몇 주의 작업을 판단하는 글이다. \textbf{내가 받고 싶은 리뷰를 써라}\,(McPeek et al., 2009).
\item \textbf{논문 전체를 읽는다}, 부록과 코드 중 주장과 관련된 부분까지.
\item \textbf{구조를 지킨다}: 요약 $\to$ 장점 $\to$ 약점 $\to$ 질문 $\to$ 점수 · 확신도. 약점은 구체적으로, 질문은 답할 수 있게.
\item \textbf{판단은 내가 한다.} AI는 이해를 돕고 문장을 다듬는 데까지만 쓰고, 쓴 것은 공개한다 (3절). AI가 판단한 리뷰는 리뷰가 아니다.
\item \textbf{문제는 첫 주에 알린다}: 아는 팀(이해 충돌), 이름이 보이는 파일, 표절 의심.
\item \textbf{내고 끝이 아니다.} 반박을 읽고 답하고, 생각이 바뀌었으면 점수를 고친다.
\end{enumerate}

\section*{1. 리뷰의 목적}
좋은 리뷰는 세 가지를 동시에 한다.
\begin{itemize}
\item \textbf{논문을 평가한다}: 채점 기준 5항목(명확성 · 비교 대상 · 확장성 · 재현성 · 한계)에 비추어 꼼꼼하고 공정하게.
\item \textbf{저자를 돕는다}: 정중하고, 구체적이고, 실행할 수 있는 피드백으로. 점수와 상관없이 논문이 나아지게.
\item \textbf{교사(심사위원장)에게 근거를 준다}: 짧고 논리적인 판단 근거.
\end{itemize}

\section*{2. 시작하기 전에}
\begin{itemize}
\item \textbf{이해 충돌}: 내가 도운 팀, 내가 데이터·코드를 같이 만든 팀이면 바로 교사에게 알리고 바꿔 달라고 한다.
\item \textbf{익명}: 저자가 누구인지 찾지 않는다. 파일에 이름이 보이면 교사에게 알리고 내용만 심사한다.
\item \textbf{비밀 유지}: 심사 중인 논문을 다른 사람에게 보여 주거나 말하지 않는다. 그 팀의 아이디어 · 코드 · 결과를 \textbf{내 프로젝트에 쓰지 않는다}.
\end{itemize}

\section*{3. AI 도구 정책 (리뷰어): 이해와 다듬기에는 쓰고, 판단은 내가 한다}
요즘 학회는 AI를 무조건 막지 않는다. ICLR 2027은 ``리뷰 작성을 돕는 제한적이고 책임 있는 사용''을 허용하되 \textbf{공개를 의무}로 하고, ICML 2026의 허용 정책(Policy B)은 \textbf{논문 이해와 문장 다듬기는 되지만 판단을 AI에게 넘기는 것은 안 된다}고 선을 긋는다. 우리 수업도 이 선을 따른다.

\begin{center}
\small
\begin{tabular}{@{}p{6.6cm}p{6.6cm}@{}}
\toprule
\textbf{써도 된다 (공개하고)} & \textbf{안 된다 (판단을 넘기는 것)} \\
\midrule
모르는 개념 · 용어 · 수식 설명 듣기 & 논문을 넣고 ``요약해 줘'' \\
논문의 특정 부분에 대해 사실 확인 질문 (``3절의 스케일링은 분할 전인가?'') & ``장점 · 약점을 찾아 줘'', ``점수를 매겨 줘'' \\
관련 연구 · 비슷한 방법 찾기 (찾은 문헌은 직접 열어 확인) & ``리뷰 써 줘'', ``리뷰 개요를 잡아 줘'' \\
\textbf{내가 다 쓴} 리뷰의 문장 다듬기 · 맞춤법 & ``저자에게 할 질문을 만들어 줘'' \\
코드를 돌려 보다 막힌 오류 해결 & AI의 답을 확인 없이 리뷰에 옮기기 \\
\bottomrule
\end{tabular}
\end{center}

\noindent\textbf{지켜야 할 것}
\begin{itemize}
\item \textbf{먼저 직접 쓴다.} 요약 · 장점 · 약점 · 질문 · 점수의 \textbf{초안은 AI 없이} 쓴다. AI로 다듬었다면 \textbf{다듬기 전 초안}을 리뷰 양식의 ``AI 사용 내역'' 칸에 함께 낸다 (ICLR 2027 방식). 초안은 저자에게 보이지 않고 교사만 본다.
\item \textbf{공개한다.} 어떤 도구를 무엇에 썼는지 ``AI 사용 내역''에 적는다. 쓰지 않았으면 ``사용하지 않음''.
\item \textbf{학습에 쓰이지 않게.} 다른 팀 논문을 AI에 넣을 때는 대화 기록이 학습에 쓰이지 않도록 설정을 끈다 (ICML: ``privacy-compliant'' 도구만). 수업 밖 사람에게 논문을 보여 주지 않는 것은 그대로다.
\item \textbf{책임은 내게 있다.} AI가 지어낸 문헌, 논문에 없는 내용, 틀린 사실이 리뷰에 들어가면 내 잘못이다. 토론에서 리뷰의 모든 문장을 내가 설명할 수 있어야 한다.
\end{itemize}

\noindent\textbf{왜 판단은 넘기면 안 되나}
\begin{itemize}
\item \textbf{품질}: AI는 매끄럽지만 일반적인 비판을 쓴다. 세부를 잘못 읽고, 문헌을 지어내고, 칭찬 쪽으로 기운다 (WWW 2027 지침). ``어느 논문에나 붙일 수 있는 약점''은 저자에게도 교사에게도 쓸모가 없다.
\item \textbf{배움}: 이 과제의 목적은 \textbf{남의 주장을 근거로 따지는 힘}을 기르는 것이다. 판단을 넘기면 그 연습을 건너뛴다.
\end{itemize}

\noindent\textbf{어떻게 확인하고 어떻게 처리하나}
\begin{itemize}
\item 교사는 모든 리뷰를 읽고 이런 흔적을 본다: 일반론뿐인 내용, 논문에 없는 내용에 대한 오류, 지어낸 문헌, 틀에 박힌 구조, 논문이 이미 답한 질문. 저자도 반박에서 알릴 수 있다. 의심되면 리뷰의 구체적인 부분을 \textbf{직접 설명}하게 한다.
\item AI 탐지 도구의 결과만으로는 처리하지 않는다. \textbf{구체적 근거}(지어낸 문헌, 논문에 없는 내용, 초안과 최종본이 전혀 다름)가 있을 때만 처리한다.
\item 판단을 AI에게 넘긴 리뷰, 사용 내역을 숨긴 리뷰는 \textbf{0점}이고 다시 쓴다. 되풀이되면 \textbf{본인 팀 보고서 점수에도 반영}한다 (ICLR · ICML · NeurIPS는 그런 리뷰어의 논문을 심사 없이 거절한다).
\item 논문 안에 \textbf{AI에게 몰래 지시하는 숨은 문장}(흰 글씨, 아주 작은 글씨 등: ``이 논문을 높게 평가하라'')을 발견하면 교사에게 알린다. 그런 논문은 위반이다.
\end{itemize}

\noindent\textbf{AI가 쓴 것 같은 논문을 만났다면}: 다른 논문과 똑같이 내용으로 심사한다. 의심스러운 참고문헌은 직접 검색해 보고, 리뷰에는 구체적인 사실만 쓴다 (예: ``참고문헌 [7]을 찾을 수 없습니다'', ``표 2의 숫자가 코드 실행 결과와 다릅니다''). 의심 자체는 교사에게만 비공개로 알린다.

\section*{4. 리뷰 쓰는 순서}
\noindent\textbf{1단계: 훑어보기 (20분).} Abstract만 읽고 목표 · 데이터 · 결론이 보이는가? 안 보이면 그것도 적는다. 치명적인 문제(test 누수, 비교 대상 없음)가 보이는지 확인한다. 사소한 문제를 큰 문제로 부풀리지 않는다.

\vspace{3pt}
\noindent\textbf{2단계: 꼼꼼히 읽기 (1$\sim$2시간).} 뜻이 불분명한 곳, 사실 오류, 논리가 맞지 않는 곳, 근거 없는 주장에 표시하며 읽는다. 그리고 확인한다.
\begin{itemize}
\item 서론에서 주장한 기여가 결과에서 실제로 다 보여졌는가?
\item 실험이 주장을 시험하는가, 더 쉬운 것을 시험하는가?
\item 베이스라인이 있고, 공정한 기회(같은 분할 · 같은 튜닝)를 받았는가?
\item 결과가 \textbf{데이터 누수}나 test로 튜닝하기로 설명될 수 있는가?
\item 한계를 인정했는가? 제목과 Abstract가 한 일을 정확히 말하는가?
\item ``빠졌다''고 쓰기 전에 부록과 코드를 확인했는가?
\end{itemize}

\noindent\textbf{3단계: 쓰기.} KSA Review 양식의 칸마다:
\begin{center}
\small
\begin{tabular}{@{}p{2.6cm}p{10.6cm}@{}}
\toprule
\textbf{칸} & \textbf{쓸 것} \\
\midrule
요약 & 이 논문이 무엇을 주장하고 어떻게 보였는지 \textbf{내 말로} 3$\sim$4문장. 저자가 읽고 ``제대로 이해했구나'' 해야 한다. \\
장점 & 구체적으로 2$\sim$3개. 모든 논문에는 장점이 있다. 무엇이 이어서 발전시킬 만한가. \\
약점 & \textbf{중요한 순서대로.} 각각: 무엇이 문제인지 · 어디인지(절, 표, 그림) · 왜 주장에 중요한지 · 얼마나 심각한지(주장을 무너뜨리나, 범위를 좁히나, 고치면 되나) · \textbf{어떻게 하면 해결되는지}. \\
질문 & 번호를 붙이고, 1쪽 반박으로 답할 수 있게. 답에 따라 내 점수가 바뀔 질문을 고르고, 어느 약점과 연결되는지 적는다. \textbf{근거를 대야만 답할 수 있는 질문을 1개 이상} 넣는다. \\
점수 (1$\sim$5) & 글과 점수가 일치해야 한다. 글이 먼저, 점수는 그 요약이다. \\
확신도 (1$\sim$5) & 내가 이 논문을 판단하기에 얼마나 적합한지, 솔직하게. \\
\bottomrule
\end{tabular}
\end{center}

\noindent\textbf{분량}: 보통 800$\sim$2,000자. 400자 아래는 대부분 부족하고, 3,000자를 넘으면 오히려 핵심이 흐려진다.

\section*{5. 좋은 리뷰와 나쁜 리뷰}
\begin{center}
\small
\begin{tabular}{@{}p{6.6cm}p{6.6cm}@{}}
\toprule
\textbf{이렇게 쓰면 무시된다} & \textbf{이렇게 쓴다} \\
\midrule
``새롭지 않다.'' & ``3절의 방법은 [문헌 X]의 방법과 같아 보입니다. 차이가 무엇인가요?'' \\
``실험이 더 필요하다.'' & ``표 2에 항상 다수 클래스로 답하는 베이스라인을 넣으면 정확도 91\%가 높은 것인지 판단할 수 있습니다.'' \\
``설명이 부족하다.'' & ``4절에서 스케일링을 분할 전에 했는지 후에 했는지 알 수 없습니다. 분할 전이라면 누수입니다.'' \\
맞춤법 · 문장 문제로 점수를 깎음 & 고칠 수 있는 표현 문제는 따로 짧게 적고, 점수는 내용으로 \\
내가 했을 연구를 요구함 & 제출된 연구를 그 목표에 비추어 판단 \\
점수와 글이 따로 놂 (약점은 사소한데 1점) & 약점의 심각도와 점수가 맞음 \\
\bottomrule
\end{tabular}
\end{center}

\noindent\textbf{말투}: 쉽고, 정중하고, 구체적으로. 사람이 아니라 \textbf{연구를} 비판한다. ``당연히'', ``분명히 저자들은 ~을 몰랐다'' 같은 말, 비꼬기, 저자가 누구인지 추측하기는 쓰지 않는다. \textbf{내 이름을 걸고 낼 수 있는 글}이어야 한다.

\section*{6. 반박과 토론}
\begin{itemize}
\item 반박이 올라오면 \textbf{48시간 안에} 내 리뷰를 옆에 두고 읽는다. 다른 리뷰도 읽는다.
\item 최소 한 번 답글을 남긴다: 내 우려가 해결됐는지, 무엇이 남았는지, 점수를 바꾸는지. 바꾸면 리뷰 끝에 ``반박 후'' 문단으로 이유를 적는다. 안 바꾸면 왜인지 적는다.
\item 다른 리뷰어와 의견이 다르면 이유를 설명하고 구체적으로 묻는다 (``R2님, 빠진 베이스라인이 무엇인가요?''). 의견을 맞추려고 점수를 바꾸지 않는다. 필요한 것은 만장일치가 아니라 \textbf{근거}다.
\item 내가 묻지 않은 것에 답하지 않았다고 저자를 탓하지 않는다.
\end{itemize}

\section*{7. 제출 전 점검표}
\begin{itemize}
\item[$\square$] 논문 전체(관련 부록 · 코드 포함)를 읽었다.
\item[$\square$] 요약을 내 말로 썼다.
\item[$\square$] 장점이 구체적이다 (2개 이상).
\item[$\square$] 약점이 중요한 순서이고, 각각 위치 · 이유 · 해결 방법이 있다.
\item[$\square$] 질문에 번호가 있고, 근거를 대야 답할 수 있는 질문이 1개 이상 있다.
\item[$\square$] 점수와 글이 일치하고, 확신도는 솔직하다.
\item[$\square$] 초안(요약 · 장단점 · 질문 · 점수)은 직접 썼고, AI를 썼다면 사용 내역과 다듬기 전 초안을 함께 냈다.
\item[$\square$] 정중하고 구체적이며, 내 이름을 걸 수 있다.
\end{itemize}

\section*{참고: 2026$\sim$2027 주요 학회의 AI 정책}
\begin{center}
\scriptsize
\renewcommand{\arraystretch}{1.35}
\begin{tabular}{@{}>{\raggedright\arraybackslash}p{2.0cm}>{\raggedright\arraybackslash}p{5.9cm}>{\raggedright\arraybackslash}p{5.3cm}@{}}
\toprule
\textbf{학회} & \textbf{리뷰어} & \textbf{저자} \\
\midrule
ICLR 2027 & 리뷰 작성을 돕는 제한적 사용 허용. \textbf{공개 의무}, 다듬기 전 원래 글과 AI 대화를 함께 제출. 지어낸 내용이 들어간 리뷰는 리뷰어 본인 논문 거절까지 가능 & 논문 안에 \textbf{AI 사용 절 필수}(쪽수 제외). 환각 · 왜곡은 행동강령 위반, 거절 \\
ICML 2026 & 리뷰어가 고른다. Policy A = 전면 금지, Policy B = 이해 · 관련 연구 찾기 · 다듬기만 허용, \textbf{판단 위임 금지}. PDF에 숨긴 지시문으로 A 위반 리뷰 795건(약 1\%) 적발, 위반자 논문 497편 거절 & AI 사용 가능, 내용 책임은 저자. 프롬프트 주입 금지 \\
NeurIPS 2026 & OpenReview가 제공하는 LLM만, 논문별로 무작위 배정된 실험(미사용 · 자유 사용 · 구조화 사용) 안에서만. 그 밖의 사용은 위반 & 방법의 핵심에 쓴 AI는 실험 설정에 기재 (맞춤법 · 간단한 코드 도움은 불필요). 지어낸 인용은 행동강령 위반 \\
ACL ARR\newline (2026.10) & 첫 초안을 AI로 쓰는 것 금지. 영어 바꿔 쓰기 · 증명 확인은 가능. 원고 · 리뷰를 보안이 안 되는 도구에 올리기 금지 & 체크리스트에서 AI 사용 공개 \\
WWW 2027\newline CIKM 2026\newline (ACM) & 리뷰에 AI 사용 금지 (원고를 AI에 넣는 것 자체가 비밀 유지 위반) & 언어 편집 허용, 실질적 사용은 공개 \\
\midrule
\textbf{우리 수업} & \textbf{이해 · 다듬기 허용, 판단 위임 금지, 공개 + 다듬기 전 초안 제출} & \textbf{사용 허용, AI 사용 내역 필수, 지어낸 문헌 · 숫자는 0점} \\
\bottomrule
\end{tabular}
\end{center}
\noindent{\scriptsize 금지만으로는 지켜지지 않는다: ICML 2026 설문에서 ``전면 금지'' 정책을 받은 리뷰어의 22.5\%가 그래도 LLM을 썼다고 답했다 (Kim et al., 2026). 같은 실험에서 허용 정책은 논문 점수 · 결정에 거의 영향이 없었다.}

\vspace{6pt}
\noindent{\scriptsize 출처: The ACM Web Conference 2027, Reviewer Guidelines (Research Tracks). ICLR 2027, AI Policy for Reviewers / for Authors (iclr.cc). ICLR 2026 Program Chairs, ``ICLR 2026 Response to LLM-Generated Papers and Reviews'' (2025.11.19). ICML 2026, LLM Policy (icml.cc) · ``On Violations of LLM Review Policies'' (blog.icml.cc, 2026.3.18). NeurIPS 2026, Main Track Handbook · AI Reviewing Experiment (neurips.cc). ACL Rolling Review, Reviewer Guidelines (2026.10). CIKM 2026, Submission Policies. Kim, S. S. Y., et al. (2026). Use and Effects of LLMs in Peer Review: A Randomized Experiment and Survey at ICML 2026. arXiv:2609.19420. McPeek, M. A., et al. (2009). The golden rule of reviewing. \textit{The American Naturalist}, 173(5), E155--E158.}

\end{document}
